Table~\ref{tab:summary} and Figures~\ref{fig:ipc}--\ref{fig:mpki} summarise the \val{n_traces_agg} traces. \prop{} ranks first of the seven predictors on every aggregate: geometric-mean IPC (\val{prop_gm_ipc}), arithmetic-mean IPC (\val{prop_mean_ipc}), pooled accuracy over all branches (\val{prop_acc}\%) and over conditional branches (\val{prop_cacc}\%), and mean MPKI (\val{prop_mpki}).

The predictors form three clear tiers. The 4 KB ChampSim predictors are far behind: bimodal (\val{bim_mpki} MPKI), gshare (\val{gsh_mpki}) and perceptron (\val{perc_mpki}). The hashed perceptron (\val{hp_mpki}) is the best ChampSim predictor and is close to the TAGE family: TAGE (\val{tage_mpki}), TAGE-SC-L (\val{tsl_mpki}) and \prop{} (\val{prop_mpki}). Relative to bimodal, the geometric-mean speedup grows from \val{gsh_speedup_bim}\% (gshare) and \val{perc_speedup_bim}\% (perceptron) to \val{hp_speedup_bim}\% (hashed perceptron), \val{tage_speedup_bim}\% (TAGE), \val{tsl_speedup_bim}\% (TAGE-SC-L) and \val{prop_speedup_bim}\% (\prop{}) (Figure~\ref{fig:speedup}). Most of the available IPC gain is therefore obtained by any of the large predictors. The differences among the top four are an order of magnitude smaller than the gap to the 4 KB predictors.

The arithmetic and geometric means give the same ranking. The arithmetic means are higher because the per-trace IPCs span more than an order of magnitude: from about 2 (648.exchange2) to well below 0.5 on the memory-bound traces (605.mcf, 620.omnetpp, 623.xalancbmk, one 602.gcc window), whose IPC is limited by cache misses whatever the predictor (Appendix~\ref{app:all}). Accuracy differences look small in percentage points because most branches are easy; MPKI expresses the same difference in the units that matter for performance. The hashed perceptron's \val{hp_acc}\% and \prop{}'s \val{prop_acc}\% correspond to a \val{vs_hp_pct_mpki}\% difference in mean MPKI.

The per-predictor means hide a very uneven distribution over traces (Figure~\ref{fig:heatmap}). With \prop{}, \val{n_prop_mpki_below_0p1} traces have an MPKI below 0.1, and \val{n_prop_mpki_above_5} have an MPKI above 5. The five traces with the highest MPKI (\val{top5_mpki_traces}) account for \val{top5_mpki_share}\% of \prop{}'s mean MPKI; all are SPECint traces (mcf, leela, xz). Mean MPKI is therefore dominated by a few hard traces, a fact that matters when interpreting the comparisons below.
